Large Language Models (LLMs) enable flexible student simulation through natural-language persona prompting, but descriptive learner profiles do not necessarily reproduce the learning-process consequences implied by process-relevant characteristics. This thesis investigates whether selected ADHD-related learning characteristics can instead be represented through explicit cognitive-process constraints that shape learner state and downstream behaviour.

A Cognitive-Process-Based (CPB) framework was developed to operationalise susceptibility to controlled distraction and sensitivity to instructional processing demand through Attention- and Working-Memory-related constraints. CPB was compared with persona-prompted student simulation using shared instructional materials, assessment questions, and a source-grounded criterion-level evaluation framework.

Three studies evaluated the approach. Study 1 verified that CPB mechanisms executed as intended and produced traceable information-state changes. Study 2 found that CPB produced clearer graded learner-condition patterns and more theory-consistent sensitivity to controlled distraction and processing demand than persona prompting. Study 3 showed that additional response-stage Language Ability and Big-Five prompts changed assessment outcomes but did not eliminate CPB's previously established constraint-related aggregate patterns.

These findings support functional process-based representation as a promising approach for modelling selected ADHD-related learning characteristics, while not claiming a complete or clinically valid simulation of ADHD cognition.
